\documentclass[10pt,a4paper]{amsart}
\usepackage{amsmath,amssymb,amsthm,graphicx}
% Retained minimal preamble for the Stacks Project chapter varieties.tex.
% Verbatim excerpt of the chapter preamble: Xy-pic setup (native preamble lines 14-18), the native theorem
% declarations with the mathematical macros (native preamble lines 154-217), and the project own no-op
% definitions of the reference/slogan/history display environments (native wrapper.tex lines 1-3) that the
% counted statement uses for its slogan. Copied mechanically; nothing added, removed or redefined.
% For commutative diagrams we use Xy-pic
\usepackage[all]{xy}

% We use 2cell for 2-commutative diagrams.
\xyoption{2cell}
% Theorem environments.
%
\theoremstyle{plain}
\newtheorem{theorem}[subsection]{Theorem}
\newtheorem{proposition}[subsection]{Proposition}
\newtheorem{lemma}[subsection]{Lemma}

\theoremstyle{definition}
\newtheorem{definition}[subsection]{Definition}
\newtheorem{example}[subsection]{Example}
\newtheorem{exercise}[subsection]{Exercise}
\newtheorem{situation}[subsection]{Situation}

\theoremstyle{remark}
\newtheorem{remark}[subsection]{Remark}
\newtheorem{remarks}[subsection]{Remarks}

\numberwithin{equation}{subsection}

% Macros
%
\def\lim{\mathop{\mathrm{lim}}\nolimits}
\def\colim{\mathop{\mathrm{colim}}\nolimits}
\def\Spec{\mathop{\mathrm{Spec}}}
\def\Hom{\mathop{\mathrm{Hom}}\nolimits}
\def\Ext{\mathop{\mathrm{Ext}}\nolimits}
\def\SheafHom{\mathop{\mathcal{H}\!\mathit{om}}\nolimits}
\def\SheafExt{\mathop{\mathcal{E}\!\mathit{xt}}\nolimits}
\def\Sch{\mathit{Sch}}
\def\Mor{\mathop{\mathrm{Mor}}\nolimits}
\def\Ob{\mathop{\mathrm{Ob}}\nolimits}
\def\Sh{\mathop{\mathit{Sh}}\nolimits}
\def\NL{\mathop{N\!L}\nolimits}
\def\CH{\mathop{\mathrm{CH}}\nolimits}
\def\proetale{{pro\text{-}\acute{e}tale}}
\def\etale{{\acute{e}tale}}
\def\QCoh{\mathit{QCoh}}
\def\Ker{\mathop{\mathrm{Ker}}}
\def\Im{\mathop{\mathrm{Im}}}
\def\Coker{\mathop{\mathrm{Coker}}}
\def\Coim{\mathop{\mathrm{Coim}}}

% Boxtimes
%
\DeclareMathSymbol{\boxtimes}{\mathbin}{AMSa}{"02}

%
% Macros for moduli stacks/spaces
%
\def\QCohstack{\mathcal{QC}\!\mathit{oh}}
\def\Cohstack{\mathcal{C}\!\mathit{oh}}
\def\Spacesstack{\mathcal{S}\!\mathit{paces}}
\def\Quotfunctor{\mathrm{Quot}}
\def\Hilbfunctor{\mathrm{Hilb}}
\def\Curvesstack{\mathcal{C}\!\mathit{urves}}
\def\Polarizedstack{\mathcal{P}\!\mathit{olarized}}
\def\Complexesstack{\mathcal{C}\!\mathit{omplexes}}
% \Pic is the operator that assigns to X its picard group, usage \Pic(X)
% \Picardstack_{X/B} denotes the Picard stack of X over B
% \Picardfunctor_{X/B} denotes the Picard functor of X over B
\def\Pic{\mathop{\mathrm{Pic}}\nolimits}
\def\Picardstack{\mathcal{P}\!\mathit{ic}}
\def\Picardfunctor{\mathrm{Pic}}
\def\Deformationcategory{\mathcal{D}\!\mathit{ef}}
\newenvironment{reference}{}{}
\newenvironment{slogan}{}{}
\newenvironment{history}{}{}
\usepackage{cleveref}
\crefname{originallemma}{original lemma}{original lemmas}
\Crefname{originallemma}{Original lemma}{Original lemmas}
\makeatletter
\newlabel{algebra-lemma-finite-type-descends}{{00QP}{source}{}{sourcepointer-1}{}}
\newlabel{algebra-lemma-finite-type-descends@cref}{{[originallemma][0][]00QP}{[1][0][]source}{}{}{}}
\makeatother
\crefname{originallemma}{original lemma}{original lemmas}
\Crefname{originallemma}{Original lemma}{Original lemmas}
\makeatletter
\newlabel{lemma-ample-after-field-extension}{{0BDC}{source}{}{sourcepointer-2}{}}
\newlabel{lemma-ample-after-field-extension@cref}{{[originallemma][0][]0BDC}{[1][0][]source}{}{}{}}
\makeatother
\crefname{originallemma}{original lemma}{original lemmas}
\Crefname{originallemma}{Original lemma}{Original lemmas}
\makeatletter
\newlabel{morphisms-lemma-finite-type-over-affine-ample-very-ample}{{01VT}{source}{}{sourcepointer-3}{}}
\newlabel{morphisms-lemma-finite-type-over-affine-ample-very-ample@cref}{{[originallemma][0][]01VT}{[1][0][]source}{}{}{}}
\makeatother
\title{Quasi-projectivity descends along an arbitrary field extension}
\author{The Stacks Project Authors}
\date{stacks-project:varieties.tex:local-snapshot-20261011; GNU FDL 1.2 or later}
\begin{document}
\maketitle
\noindent\textit{Editorial note.} All delivered mathematical text is the author's own, copied from the pinned Stacks Project chapter file varieties.tex; no statement or proof step was rewritten, reordered or removed.\par\medskip
\noindent\textit{Editorial note.} The proof is complete as written and ends with the quasi-projectivity conclusion; no step is referred to the reader.\par\medskip
\noindent\textit{Editorial note.} The counted spans carry 4 navigation stubs and no bibliography citation.\par\medskip
\noindent\textit{Editorial note.} The preamble is the chapter's retained minimal preamble (sources/stacks-varieties/retained-preamble.tex lines 1-77), a verbatim excerpt of the native preamble.\par\medskip
\noindent\textit{Editorial note.} The label lemma-quasi-projective-after-field-extension is defined only in varieties.tex, so no cross-chapter restatement had to be weighed for this card.\par\medskip
\noindent\textit{Allowed standard prerequisites (editorial).}\begin{itemize}
\item Finiteness of generation descends along a field extension. Algebra, Lemma algebra-lemma-finite-type-descends, invoked at proof lines 2830-2831. Supplies that the coordinate rings of affine opens of X are finitely generated k-algebras.
\item Ampleness descends along a field extension. Same chapter, Lemma lemma-ample-after-field-extension, invoked at proof lines 2838-2839. Supplies the ample invertible sheaf on X.
\item Quasi-projectivity is equivalent to local finite type, quasi-compactness and an ample invertible sheaf. Morphisms, Lemma morphisms-lemma-finite-type-over-affine-ample-very-ample, invoked at proof lines 2835-2837 and again at 2840-2841. Supplies the unpacking of the definition used at the start and the closing step.
\end{itemize}\medskip
\ifdefined\hypertarget\hypertarget{sourcepointer-1}{}\fi
\noindent\textit{Source reference.} Original lemma 00QP: sources/stacks-chow/tags.txt, source lines 975--975. Pointer outside the counted span, invoked once for context; the load-bearing argument of the counted statement is the author's proof reproduced in full in the proof span.\par\medskip
\ifdefined\hypertarget\hypertarget{sourcepointer-2}{}\fi
\noindent\textit{Source reference.} Original lemma 0BDC: sources/stacks-varieties/pinned-original-source.tex, source lines 2740--2740. Pointer outside the counted span, invoked once for context; the load-bearing argument of the counted statement is the author's proof reproduced in full in the proof span.\par\medskip
\ifdefined\hypertarget\hypertarget{sourcepointer-3}{}\fi
\noindent\textit{Source reference.} Original lemma 01VT: sources/stacks-chow/tags.txt, source lines 2458--2458. Pointer outside the counted span, invoked once for context; the load-bearing argument of the counted statement is the author's proof reproduced in full in the proof span.\par\medskip
\begin{lemma}
\label{lemma-quasi-projective-after-field-extension}
Let $k$ be a field. Let $X$ be a scheme over $k$. If $X_K$ is quasi-projective
over $K$ for some field extension $K/k$, then $X$ is quasi-projective
over $k$.
\end{lemma}

\begin{proof}
By definition a morphism of schemes $g : Y \to T$ is quasi-projective
if it is locally of finite type, quasi-compact, and there exists
a $g$-ample invertible sheaf on $Y$.
Let $K/k$ be a field extension such that $X_K$ is quasi-projective
over $K$. Let $\Spec(A) \subset X$ be an affine open. Then $U_K$ is an
affine open subscheme of $X_K$, hence $A_K$ is a $K$-algebra of finite type.
Then $A$ is a $k$-algebra of finite type by
Algebra, Lemma \ref{algebra-lemma-finite-type-descends}.
Hence $X \to \Spec(k)$ is locally of finite type.
Since $X_K \to \Spec(K)$ is quasi-compact, we see that $X_K$ is
quasi-compact, hence $X$ is quasi-compact, hence $X \to \Spec(k)$
is of finite type. By Morphisms, Lemma
\ref{morphisms-lemma-finite-type-over-affine-ample-very-ample}
we see that $X_K$ has an ample invertible sheaf.
Then $X$ has an ample invertible sheaf by
Lemma \ref{lemma-ample-after-field-extension}.
Hence $X \to \Spec(k)$ is quasi-projective by Morphisms, Lemma
\ref{morphisms-lemma-finite-type-over-affine-ample-very-ample}.
\end{proof}

\end{document}
